\section{Del dato tabular al contexto formal}
\label{sec:contexto}

El Análisis Formal de Conceptos trabaja sobre tablas binarias: cada objeto tiene o no tiene cada atributo. Los datos reales traen columnas numéricas (radio, textura, área\dots) y categóricas (diagnóstico). Esta sección formaliza los tres pasos que llevan de una tabla así a un contexto formal: la discretización de las variables numéricas (\S\ref{sec:discretizacion} y \S\ref{sec:numero-intervalos}), el escalado conceptual (\S\ref{sec:escalado} y \S\ref{sec:propiedades-escalado}) y la matriz de incidencia (\S\ref{sec:contextos}). En el código, \path{cuant_type.py} y \path{quantiles.py} discretizan, y \path{scaler.py} escala y produce la matriz; todos están en \path{src/preprocessing/}, y su documentación en \path{docs/codigo/preprocessing/}.


\subsection{Contextos formales y matriz de incidencia}
\label{sec:contextos}

\begin{definicion}[contexto formal]\label{def:contexto}
Un \emph{contexto formal} es una terna $\K = (G, M, I)$ formada por un conjunto finito~$G$ de \emph{objetos}, un conjunto finito~$M$ de \emph{atributos} y una relación $I \subseteq G \times M$. Si $(g, m) \in I$ se escribe $g \mathrel{I} m$: \enquote{el objeto~$g$ tiene el atributo~$m$}.
\end{definicion}

\begin{definicion}[fila y columna]\label{def:fila-columna}
Para $g \in G$ y $m \in M$, $g' = \{ m \in M : g \mathrel{I} m \}$ es la \emph{fila} de~$g$ y $m' = \{ g \in G : g \mathrel{I} m \}$ es la \emph{columna} de~$m$.
\end{definicion}

\begin{definicion}[matriz de incidencia]\label{def:matriz}
Con $G = \{g_1, \dots, g_n\}$ y $M = \{m_1, \dots, m_p\}$ en un orden fijo, la \emph{matriz de incidencia} es $X \in \{0, 1\}^{n \times p}$ con $X_{ij} = 1$ si y solo si $g_i \mathrel{I} m_j$. Se dibuja como una \emph{tabla de cruces}.
\end{definicion}

\begin{observacion}[representación]\label{obs:representacion}
Fijados los órdenes, la matriz determina el contexto. Se puede guardar de forma explícita, como un \texttt{DataFrame} booleano, o implícita, con cada fila~$g'$ como un entero usado como vector de bits; esta es la que usa el código de la Sección~\ref{sec:conceptos}.
\end{observacion}


\subsection{Contextos multivaluados y escalado conceptual}
\label{sec:escalado}

\begin{definicion}[contexto multivaluado]\label{def:multivaluado}
Un \emph{contexto multivaluado} $(G, M, W, I)$ tiene objetos~$G$, atributos multivaluados~$M$, valores~$W$ y una relación $I \subseteq G \times M \times W$ en la que cada par $(g, m)$ tiene a lo sumo un valor, que se escribe $m(g)$; si no tiene ninguno, el valor \emph{falta}. $D_m$ es el conjunto de objetos en los que $m$ está definido y $V_m = \{ m(g) : g \in D_m \}$ el de sus \emph{valores observados}.
\end{definicion}

\begin{observacion}[modelo relacional]\label{obs:relacional}
Un contexto multivaluado es una tabla de una base de datos: los objetos son las tuplas, los atributos son las columnas y un valor faltante es un \texttt{NULL}.
\end{observacion}

\begin{definicion}[escala conceptual]\label{def:escala}
Una \emph{escala} para~$m$ es un contexto formal $\mathbb{S}_m = (S, N_m, I_m)$ cuyos objetos incluyen los valores observados de~$m$; sus atributos~$N_m$ son los atributos binarios con los que se describe cada valor. En el código, $S = V_m$.
\end{definicion}

\begin{definicion}[contexto derivado]\label{def:derivado}
El \emph{contexto derivado} de $(G, M, W, I)$ con escalas~$\mathbb{S}_m$ es el contexto formal $(G, N, J)$ con
\[
  N = \bigcup_{m \in M} \{m\} \times N_m, \qquad g \mathrel{J} (m, a) \iff g \in D_m \text{ y } m(g) \mathrel{I_m} a.
\]
Cada columna~$m$ se reemplaza por un bloque de columnas binarias; si a~$g$ le falta el valor de~$m$, su fila en ese bloque queda vacía.
\end{definicion}

\begin{definicion}[escalas nominal y ordinal]\label{def:escalas}
Para valores $v_1 < \dots < v_r$, la \emph{escala nominal} es $(V, V, =)$, la \emph{ordinal creciente} es $(V, V, \le)$ y la \emph{ordinal decreciente} es $(V, V, \ge)$. En el contexto derivado, sus atributos se escriben $(m{=}w)$, $(m{\le}w)$ y $(m{\ge}w)$, y son las columnas \texttt{m=w}, \texttt{m<=w} y \texttt{m>=w} de la función \texttt{scaling} de \path{scaler.py}.
\end{definicion}

\begin{ejemplo}[tres valores]\label{ej:escalas}
Con $V = \{0, 1, 2\}$, la escala nominal es la matriz identidad, y la ordinal creciente es triangular superior: el valor~0 tiene $(\le 0)$, $(\le 1)$ y $(\le 2)$; el valor~1, $(\le 1)$ y $(\le 2)$; y el valor~2, solo $(\le 2)$.
\end{ejemplo}

\begin{observacion}[escala interordinal]\label{obs:interordinal}
Si se usan a la vez las dos escalas ordinales (la escala \emph{interordinal} de \textcite{ganter1999}), los atributos $(m{\ge}v)$ y $(m{\le}w)$ describen juntos el intervalo $v \le m(g) \le w$. Se obtiene concatenando las salidas de \texttt{scaling} con \texttt{"le"} y con \texttt{"ge"}.
\end{observacion}


\subsection{Propiedades del escalado}
\label{sec:propiedades-escalado}

Sea $m$ un atributo con valores observados $v_1 < \dots < v_r$, escalado sobre~$V_m$ como en el código. Las propiedades siguientes salen directamente de las definiciones, usando que $m(g)$ es único.

\begin{proposicion}[la escala nominal induce una partición]\label{prop:particion}
Los conjuntos $(m{=}v)'$, con $v \in V_m$, son no vacíos, disjuntos dos a dos y su unión es~$D_m$. Sin valores faltantes, cada fila tiene exactamente un~1 en el bloque de~$m$.
\end{proposicion}

\begin{proposicion}[la escala ordinal induce una cadena]\label{prop:cadena}
$(m{\le}v_1)' \subseteq \dots \subseteq (m{\le}v_r)' = D_m$ y $D_m = (m{\ge}v_1)' \supseteq \dots \supseteq (m{\ge}v_r)'$.
\end{proposicion}

\begin{corolario}[atributo universal]\label{cor:universal}
Sin valores faltantes, todos los objetos tienen $(m{\le}v_r)$ y $(m{\ge}v_1)$: sus columnas están llenas. Un atributo así no cambia el retículo y se puede quitar (Proposición~\ref{prop:universal}).
\end{corolario}

\begin{observacion}[regularidades de la escala]\label{obs:fondo}
Por las Proposiciones \ref{prop:particion} y~\ref{prop:cadena}, todo objeto con $(m{\le}v_i)$ tiene $(m{\le}v_j)$ si $i < j$, y ninguno tiene a la vez $(m{=}v)$ y $(m{=}w)$ con $v \ne w$. Son regularidades de la escala, no de los datos; la \S\ref{sec:fondo} las separa de las implicaciones que sí vienen de los datos.
\end{observacion}

\begin{proposicion}[el escalado no pierde información]\label{prop:sin-perdida}
Para $g \in D_m$, el valor $m(g)$ se recupera de la fila~$g'$: es el único~$v$ con $(m{=}v) \in g'$, el menor~$v$ con $(m{\le}v) \in g'$ o el mayor~$v$ con $(m{\ge}v) \in g'$, según la escala.
\end{proposicion}

\begin{corolario}\label{cor:filas}
Dos objetos sin valores faltantes tienen la misma fila en el contexto derivado si y solo si tienen los mismos valores en todas las variables. El contexto derivado tiene $\sum_{m} \lvert V_m \rvert$ atributos.
\end{corolario}

Así, toda la pérdida de información del pipeline ocurre en la discretización. Con las 30~variables de Breast Cancer Wisconsin y $k = 11$ intervalos por variable, el contexto tiene a lo sumo 330~atributos.


\subsection{Discretización}
\label{sec:discretizacion}

Una variable continua toma casi tantos valores distintos como objetos hay, y escalarla directamente daría casi una columna por objeto. La discretización agrupa los valores en pocos intervalos ordenados.

\begin{definicion}[discretización por puntos de corte]\label{def:discretizacion}
Dados cortes $c_0 < c_1 < \dots < c_k$ con $c_0 \le \min x$ y $c_k \ge \max x$, los intervalos son $I_0 = [c_0, c_1]$ e $I_j = (c_j, c_{j+1}]$ para $1 \le j < k$, y la discretización asigna a cada objeto el índice $\delta(g) = j$ del intervalo que contiene a~$x(g)$. Es la convención de pandas.
\end{definicion}

\begin{proposicion}[la discretización conserva el orden]\label{prop:orden}
Si $x(g) \le x(h)$, entonces $\delta(g) \le \delta(h)$.
\end{proposicion}

Por eso se puede aplicar una escala ordinal a los códigos, y cada atributo ordinal se lee en la variable original: $\delta(g) \le j$ si y solo si $x(g) \le c_{j+1}$. Con los cortes del radio de la \S\ref{sec:ejemplo}, por ejemplo, \texttt{radio<=1} significa \enquote{radio${}\le 17.13$}. La discretización es el único paso con pérdida: los objetos cuyos valores caen en un mismo intervalo se vuelven indistinguibles en esa variable. El proyecto ofrece tres formas de elegir los cortes.

\begin{definicion}[igual ancho]\label{def:ancho}
$c_j = a + j(b - a)/k$, con $a = \min x$ y $b = \max x$. Todos los intervalos miden lo mismo, pero un valor atípico estira el rango y puede dejar casi todos los objetos en uno o dos intervalos.
\end{definicion}

\begin{definicion}[igual frecuencia]\label{def:frecuencia}
Los cortes son los cuantiles empíricos $c_j = \hat{Q}(j/k)$, con la interpolación lineal que usan por defecto NumPy y pandas, el tipo~7 de \textcite{hyndman1996}. Con $k = 4$, los cortes interiores son los cuartiles.
\end{definicion}

\begin{proposicion}[intervalos balanceados]\label{prop:balance}
Si los $n$~valores son distintos, cada intervalo de igual frecuencia contiene, en aritmética exacta, $\lfloor n/k \rfloor$ o $\lceil n/k \rceil$ objetos.
\end{proposicion}

\begin{observacion}[empates]\label{obs:empates}
Con valores repetidos, varios cuantiles pueden coincidir; la implementación elimina los cortes repetidos, así que pueden quedar menos de $k$~intervalos. Por redondeo, algún intervalo puede diferir en un objeto del tamaño de la Proposición~\ref{prop:balance}.
\end{observacion}

\begin{definicion}[$k$-medias en una dimensión]\label{def:kmedias}
Se buscan grupos $C_1, \dots, C_k$ y centroides $\mu_1, \dots, \mu_k$ que minimicen $\SSE(C, \mu) = \sum_{i} \sum_{x \in C_i} (x - \mu_i)^2$. El \emph{algoritmo de Lloyd} \parencite{lloyd1982} alterna dos pasos hasta que los centroides no cambian: asignar cada valor al centroide más cercano y reemplazar cada centroide por la media de su grupo.
\end{definicion}

\begin{lema}[la media minimiza la suma de cuadrados]\label{lem:media}
Si $S$ es un multiconjunto finito de reales con media~$\bar{s}$, para todo $c \in \R$ se tiene $\sum_{x \in S} (x - c)^2 = \sum_{x \in S} (x - \bar{s})^2 + \lvert S \rvert \, (\bar{s} - c)^2$.
\end{lema}

\begin{teorema}[convergencia del algoritmo de Lloyd]\label{teo:lloyd}
La SSE no aumenta en ningún paso, disminuye estrictamente en cada iteración en que cambia algún centroide, y los centroides dejan de cambiar después de un número finito de iteraciones.
\end{teorema}

La asignación no aumenta ningún término de la SSE, la actualización tampoco, por el Lema~\ref{lem:media}, y como hay finitas formas de repartir los valores, la SSE no puede bajar indefinidamente.

\begin{observacion}[óptimo local]\label{obs:optimo-local}
El resultado es un óptimo local, que depende de los centroides iniciales; la implementación los pone en percentiles fijos, así que es reproducible. En una dimensión, el óptimo global se puede calcular con programación dinámica \parencite{fisher1958, wang2011}.
\end{observacion}

\begin{proposicion}[regiones de asignación]\label{prop:voronoi}
Si $\mu_1 < \dots < \mu_k$, el centroide más cercano a~$x$ es~$\mu_j$ si y solo si $x$ está entre los puntos medios $(\mu_{j-1} + \mu_j)/2$ y $(\mu_j + \mu_{j+1})/2$.
\end{proposicion}

Por eso \texttt{kmeans\_1d}, al terminar, corta en los puntos medios entre centroides consecutivos, y los códigos conservan el orden (Proposición~\ref{prop:orden}).


\subsection{Número de intervalos}
\label{sec:numero-intervalos}

Con pocos intervalos se pierden distinciones entre objetos (Corolario~\ref{cor:filas}); con muchos, el contexto tiene más atributos y el número de conceptos puede crecer mucho, que es el principal reto del dataset de Spotify. Las reglas siguientes vienen de la teoría de histogramas.

\begin{regla}[\cite{sturges1926}]\label{regla:sturges}
$k = \lceil \log_2 n \rceil + 1$. Sturges idealizó el histograma de una variable normal como una binomial con $k$~intervalos, cuyas frecuencias suman $2^{k-1} = n$. Con $n = 569$ da $k = 11$.
\end{regla}

\begin{teorema}[\cite{scott1979}]\label{teo:amise}
Si la densidad~$f$ es suave y $R(f') = \int f'(x)^2 \, dx$ es finito, el error cuadrático medio integrado de un histograma de ancho~$h$ con $n$~datos es, asintóticamente, $\AMISE(h) = 1/(nh) + h^2 R(f')/12$.
\end{teorema}

\begin{corolario}[ancho óptimo]\label{cor:ancho-optimo}
La AMISE se minimiza en $h^{*} = \bigl( 6 / (n R(f')) \bigr)^{1/3}$, que decrece como $n^{-1/3}$; la regla de Sturges, en cambio, crece solo como $\log_2 n$.
\end{corolario}

\begin{corolario}[regla de Scott]\label{cor:scott}
Para una normal con desviación~$\sigma$, $R(f') = 1/(4\sqrt{\pi}\,\sigma^3)$ y $h^{*} = (24\sqrt{\pi})^{1/3} \sigma\, n^{-1/3} \approx 3.49\, \sigma\, n^{-1/3}$. La implementación estima~$\sigma$ con la desviación estándar muestral.
\end{corolario}

\begin{regla}[\cite{freedman1981}]\label{regla:fd}
$h = 2 \IQR\, n^{-1/3}$, con $\IQR$ el rango intercuartílico. Es la variante robusta de la regla de Scott: el IQR, a diferencia de la desviación estándar, no se ve afectado por los valores atípicos.
\end{regla}

\begin{observacion}[del ancho al número de intervalos]\label{obs:ancho-k}
Con un ancho~$h$, $k = \lceil (\max x - \min x)/h \rceil$; si la dispersión es nula se usa la regla de Sturges, y \texttt{discretize} obliga a que $k \ge 2$. Las reglas se diseñaron para intervalos de igual ancho: usarlas con igual frecuencia o con $k$-medias es una heurística.
\end{observacion}


\subsection{Ejemplo completo}
\label{sec:ejemplo}

Se toman seis objetos con dos variables numéricas y una categórica (valores ilustrativos, inspirados en Breast Cancer Wisconsin). La tabla de la izquierda tiene los valores originales y la de la derecha sus códigos después de discretizar.

\begin{center}
\small
\begin{tabular}{cccc}
  \toprule
  Objeto & radio & textura & diagnóstico \\
  \midrule
  $g_1$ & 12.0 & 15.7 & B \\
  $g_2$ & 13.5 & 21.3 & B \\
  $g_3$ & 20.1 & 17.8 & M \\
  $g_4$ & 11.2 & 25.0 & B \\
  $g_5$ & 17.8 & 19.9 & M \\
  $g_6$ & 14.6 & 14.1 & B \\
  \bottomrule
\end{tabular}
\qquad\qquad
\begin{tabular}{cccc}
  \toprule
  Objeto & radio & textura & diagnóstico \\
  \midrule
  $g_1$ & 0 & 0 & B \\
  $g_2$ & 0 & 1 & B \\
  $g_3$ & 2 & 1 & M \\
  $g_4$ & 0 & 2 & B \\
  $g_5$ & 2 & 1 & M \\
  $g_6$ & 1 & 0 & B \\
  \bottomrule
\end{tabular}
\end{center}

Con intervalos de igual ancho y $k = 3$, los cortes del radio son 11.2, 14.17, 17.13 y~20.1 (pandas baja el primero a 11.19 para incluir el mínimo), y los de la textura 14.1, 17.73, 21.37 y~25.0. Con la escala ordinal creciente para el radio y la textura, y la nominal para el diagnóstico, el contexto tiene $\lvert G \rvert = 6$ objetos y $\lvert M \rvert = 8$ atributos:

\begin{center}
\small
\begin{tabular}{c|ccc|ccc|cc}
  & \multicolumn{3}{c|}{radio} & \multicolumn{3}{c|}{textura} & \multicolumn{2}{c}{diagnóstico} \\
  & $\le 0$ & $\le 1$ & $\le 2$ & $\le 0$ & $\le 1$ & $\le 2$ & $=\mathrm{B}$ & $=\mathrm{M}$ \\ \hline
  $g_1$ & $\times$ & $\times$ & $\times$ & $\times$ & $\times$ & $\times$ & $\times$ &          \\
  $g_2$ & $\times$ & $\times$ & $\times$ &          & $\times$ & $\times$ & $\times$ &          \\
  $g_3$ &          &          & $\times$ &          & $\times$ & $\times$ &          & $\times$ \\
  $g_4$ & $\times$ & $\times$ & $\times$ &          &          & $\times$ & $\times$ &          \\
  $g_5$ &          &          & $\times$ &          & $\times$ & $\times$ &          & $\times$ \\
  $g_6$ &          & $\times$ & $\times$ & $\times$ & $\times$ & $\times$ & $\times$ &
\end{tabular}
\end{center}

Las columnas \enquote{radio${}\le 2$} y \enquote{textura${}\le 2$} están llenas (Corolario~\ref{cor:universal}); \enquote{radio${}\le 0$} está contenida en \enquote{radio${}\le 1$} (Proposición~\ref{prop:cadena}); las dos columnas del diagnóstico parten los objetos en $\{g_1, g_2, g_4, g_6\}$ y $\{g_3, g_5\}$ (Proposición~\ref{prop:particion}); y $g_3$ y~$g_5$ tienen la misma fila aunque sus valores originales son distintos, que es la pérdida de información de la discretización (Corolario~\ref{cor:filas}). El código que produce este ejemplo está en \path{docs/codigo/preprocessing/scaler.md}.
